%==================================================================FRONT PAGE AND TOC
% For article only
\mode<presentation:0>{\thispagestyle{empty}\maketitle}

% For presentation only
\mode<presentation| article:0| handout:0>{
    \begin{frame}<article:0>[label=portada]
    \titlepage
    \end{frame}%Fin del frame
}

% For handout only
\mode<handout>{
  \begin{frame}[label=portada]
    \maketitle
  \end{frame}
}

% %% TABLE OF CONTENTS
% \begin{frame}[label=toc]
%     \mode<article:0>{\frametitle{Contents}}
%     \mode<presentation>{\small}
%     \tableofcontents[hidesubsections]
% \end{frame}

\begin{frame}[label=portada]
 \titlepage
\end{frame}
\begin{frame}[label=toc]
\tableofcontents[hidesubsections]
\end{frame}

%%==================================================================S INTRODUCTION
\section{Herramientas}
%%==================================================================F 
\begin{frame}
    \frametitle{Herramientas}
    \begin{enumerate}
        \item GRASS LiDAR tools 
	\begin{itemize}
            \item \texttt{v.outlier}
            \item \texttt{v.lidar.edgedetetion}
            \item \texttt{v.lidar.growing}
            \item \texttt{v.lidar.correction}
            \item \texttt{v.surf.bspline}
	\end{itemize}
        \item LibLAS
	\begin{itemize}
            \item \texttt{las2las}, \texttt{las2txt}, \texttt{txt2las}, \texttt{lasdiff}, \texttt{lasinfo} 
	\end{itemize}
        \item LAStools
	\begin{itemize}
            \item \texttt{las2las}, \texttt{las2txt}, \texttt{txt2las},
                \texttt{shp2las}, \texttt{lasground}, \texttt{las2dem}, \texttt{txt2tin}
	\end{itemize}
        \item PCL (Point Cloud Library)
        \item PDAL (Point Data Abstraction Library)
    \end{enumerate}
\end{frame}
%%==================================================================S
\section{Análisis LiDAR en GRASS}
%%==================================================================Sb
\begin{frame}
    \frametitle{Geographic Resources Analysis Support System\hfill
        \includegraphics[width=0.7cm]{images/grasslogo_transp_big.png}} 
     \begin{itemize}
        \item Conocido como \alert<1>{GRASS}: GIS para manipulación y
            análisis de datos geoespaciales.
        \item Como es un programa de \alert{códico abierto} con un montón de
            librerias, es muy adecuado para la implementación de nuestras
            propias herramientas.
        \item Escrito en C/C++, lo que permite implementar programas
            computacionalmente pesados $\Rightarrow$ \alert{ideal} para trabajar
            con LiDAR.
        \item GRASS $\geq$ 6 incluye indexación topológica $\Rightarrow$
            \alert{fatal} para trabajar con LiDAR.  \item Hay versiones para
            GNU/Linux, MS-Windows, Mac-OSX, SUN, etc.
            \alert{\textexclamdown\textexclamdown Todos lo pueden utilizar!!}
        \item Más información en:
            \beamergotobutton{\url{http://grass.osgeo.ogr}}
    \end{itemize}
\end{frame}
%%==================================================================Sb
\subsection{Importando nube de puntos}
%%==================================================================F 
\begin{frame}
 \frametitle{Importar como raster: \LARGE\path{r.in.xyz}}
\end{frame}
%%==================================================================F 
\begin{frame}
 \frametitle{Importar como vectorial: \LARGE\path{v.in.ascii}}
\end{frame}
%%==================================================================Sb
\subsection{Eliminando puntos groseros}
%%==================================================================F 
\begin{frame}[fragile,shrink=5]
  \frametitle{\path{v.outlier}}
  \begin{beamerboxesrounded}[shadow=true]{\textbf{\path{v.lidar.edgedetection}}
    \texttt{ input=name output=name outlier=name [soe=value] [son=value] 
    [lambda\_i=value] [thres\_o=value]}}
    \begin{itemize}
      \item \textbf{input}: name of the input vector map
      \item \textbf{output}: name of the output vector map
      \item \textbf{soe}: Interpolation spline step value in e-w direction
      \item \textbf{son}: Interpolation spline step value in n-s direction
      \item \textbf{lambda\_i}: Thychonov regularization weigth. (default: 0.1)
      \item \textbf{thres\_o}: Threshold for the outliers. (default: 50)
    \end{itemize}
  \end{beamerboxesrounded}
  \pause
  \begin{beamerboxesrounded}[shadow=true]{Sintaxis}
\scriptsize
%\testcode
\begin{verbatim}
GRASS6 (stuttgart): > # Limpiamos el primer impulso
GRASS6 (stuttgart): > v.outlier in=stuttgart_first out=station_first \
> outlier=station_first_out soe=10 son=10 thres_o=25
GRASS6 (stuttgart): > # Limpiamos el ultimo impulso
GRASS6 (stuttgart): > v.outlier in=stuttgart_last out=station_last \
> outlier=station_last_out soe=10 son=10 thres_o=30
\end{verbatim}

\end{beamerboxesrounded}
\end{frame}
%%==================================================================F 
\begin{frame}
 \frametitle{Visualización de \LARGE\path{v.outlier}}
\end{frame}
%%==================================================================Sb
\subsection{Detección de borde de objetos}
%%==================================================================F 
\begin{frame}[fragile,shrink=10]
  \frametitle{\path{v.lidar.edgedetection}}
  \begin{beamerboxesrounded}[shadow=true]{\textbf{\path{v.lidar.edgedetection}}
    \texttt{ input=name output=name first=name see=value sen=value 
    [lambda\_g=value] [tgh=value] [tgl=value] [theta\_g=value] [lambda\_r=value]}}
   \begin{itemize}
    \item \textbf{input}: name of the input vector map
    \item \textbf{output}: name of the output vector map
    \item \textbf{see}: Interpolation spline step value in e-w direction
    \item \textbf{sen}: Interpolation spline step value in n-s direction
    \item \textbf{lambda\_g}: Regularization weight in gradient evaluation (default: 0.01)
    \item \textbf{tgh}: High gradient threshold for edge classification (default: 6)
    \item \textbf{tgl}: Low gradient threshold for edge classification (default: 3)
    \item \textbf{theta\_g}: Angle range for same direction detection (default: 0.26)
    \item \textbf{lambda\_r}: Regularization weight in residual evaluation (default: 2)
   \end{itemize}
  \end{beamerboxesrounded}
 \pause
 \begin{beamerboxesrounded}[shadow=true]{Sintaxis}
\scriptsize
%\testcode
\begin{verbatim}
GRASS6 (stuttgart): > v.lidar.edgedetection in=station_last see=4 sen=4 \
> lambda_g=0.1 out=station_edge
\end{verbatim}
\end{beamerboxesrounded}
\end{frame}
%%==================================================================F 
\pgfdeclareimage[width=0.65\textwidth]{edge}{images/edge}
\begin{frame}
 \frametitle{Visualización de \LARGE\path{v.lidar.edgedetection}}
\begin{columns}
  \begin{column}{0.7\textwidth}
    \begin{center}
	\begin{tikzpicture}
    	  \pgftext[bottom,left,at={\pgfpointxy{0}{0}}]{\pgfuseimage{edge}}
	\end{tikzpicture}
   \end{center}
  \end{column}
  \begin{column}{0.25\textwidth}
	\begin{enumerate}
	 \item \textcolor{red}{Borde}
	 \item \textcolor{yellow!95!black}{Terreno}
	\end{enumerate}
  \end{column}
\end{columns}
\end{frame}
%%==================================================================Sb
\subsection{Relleno de los límites de los objetos}
%%==================================================================F 
\begin{frame}[fragile,shrink=10]
  \frametitle{\path{v.lidar.growing}}
  \begin{beamerboxesrounded}[shadow=true]{\textbf{\path{v.lidar.growing}}
    \texttt{ input=name output=name first=name [tj=value] [td=value]}}
    \begin{itemize}
     \item \textbf{input}: name of the input vector map
     \item \textbf{output}: name of the output vector map
     \item \textbf{first}: name of the input vector map of the first pulse points
     \item \textbf{tj}: Threshold in planimetric units for considering two 
         measurements as corresponding first and double pulse. (default 0.2)
     \item td: Threshold in height units for classifying a point as double pulse. 
         (default: 0.6)
    \end{itemize}
  \end{beamerboxesrounded}
  \pause
  \begin{beamerboxesrounded}[shadow=true]{Sintaxis}
\scriptsize
\begin{verbatim}
GRASS6 (stuttgart): > # Se disminuye primero la resolucion
GRASS6 (stuttgart): > g.region -p res=2
GRASS6 (stuttgart): > v.lidar.growing in=station_edge out=station_grow \
> first=station_first
\end{verbatim}
\end{beamerboxesrounded}
\begin{itemize}
 \item \alert{Importante!!} Bajar la resolución de la región
\end{itemize}
\end{frame}
%%==================================================================F 
\pgfdeclareimage[width=0.65\textwidth]{grow}{images/grow}
\begin{frame}
 \frametitle{Visualización de \LARGE\path{v.lidar.growing}}
\begin{columns}
  \begin{column}{0.65\textwidth}
    \begin{center}
	\begin{tikzpicture}
    	  \pgftext[bottom,left,at={\pgfpointxy{0}{0}}]{\pgfuseimage{grow}}
	\end{tikzpicture}
   \end{center}
  \end{column}
  \begin{column}{0.35\textwidth}
	\begin{enumerate}
	 \item \textcolor{yellow!95!black}{Terreno}
	 \item \textcolor{green}{Terreno doble eco}
	 \item \textcolor{blue!90!black}{Objeto doble eco}
	 \item \textcolor{red}{Objeto}
	\end{enumerate}
  \end{column}
\end{columns}
\end{frame}
%%==================================================================Sb
\subsection{Corrección final}
%%==================================================================F 
\begin{frame}[fragile,shrink=5]
  \frametitle{\path{v.lidar.correction}}
  \begin{beamerboxesrounded}[shadow=true]{\textbf{\path{v.lidar.correction}}
    \texttt{ input=name  output=name  terrain=name  sce=value  scn=value 
    [lambda\_c=value]  [tch=value]  [tcl=value]}}
    \begin{itemize}
     \item \textbf{input}: name of the input vector map
     \item \textbf{output}: name of the output vector map
     \item \textbf{terrain}: name of the terrain output vector map
     \item \textbf{sce}: Interpolation spline step value in e-w direction
     \item \textbf{scn}: Interpolation spline step value in n-s direction
     \item \textbf{lambda\_c}: Regularization weight in reclassification evaluation. (default 1)
     \item \textbf{tch}: High difference threshold (terrain/object). (default 2)
     \item \textbf{tcl}: Low differnce threshold (object/terrain). (default 1)
    \end{itemize}
  \end{beamerboxesrounded}
  \pause
  \begin{beamerboxesrounded}[shadow=true]{Sintaxis}
\scriptsize
\begin{verbatim}
GRASS6 (stuttgart): > v.lidar.correction in=station_grow sce=60 scn=60 \
> lambda_c=2 tcl=0.1 out=station_corr terrain=station_terr
\end{verbatim}
  \end{beamerboxesrounded}
\end{frame}
%%==================================================================F 
\pgfdeclareimage[width=0.65\textwidth]{corr}{images/correction}
\begin{frame}
 \frametitle{Visualización de \LARGE\path{v.lidar.correction}}
\begin{columns}
  \begin{column}{0.65\textwidth}
    \begin{center}
	\begin{tikzpicture}
    	  \pgftext[bottom,left,at={\pgfpointxy{0}{0}}]{\pgfuseimage{corr}}
	\end{tikzpicture}
   \end{center}
  \end{column}
  \begin{column}{0.35\textwidth}
	\begin{enumerate}
	 \item \textcolor{yellow!95!black}{Terreno}
	 \item \textcolor{green}{Terreno doble eco}
	 \item \textcolor{blue!90!black}{Objeto doble eco}
	 \item \textcolor{red}{Objeto}
	\end{enumerate}
  \end{column}
\end{columns}
\end{frame}






%%==================================================================S INTRODUCTION
\section{¿Qué es MATLAB?}
%%==================================================================F 
\begin{frame}[label=que_es]
    \frametitle{¿Qué es MATLAB?}
    \begin{enumerate}
        \item \alert{MAT}rix \alert{LAB}oratory
        \item Propiedad de MathWorks (\url{http://www.mathworks.es/})
        \item Programa para cálculo científico y de ingeniería
	\begin{itemize}
	 \item Extensible (\emph{toolboxes})
         \item Lenguaje de programación proprio interpretado (\alert{M})
         \item Visualización interactiva de gráficos
         \item Integrable con aplicaciones escritas en otros lenguajes
	\end{itemize}
        \item Existen implementaciones libres del lenguaje M
	\begin{itemize}
	 \item Scilab (\url{http://www.scilab.org})
	 \item Octave (\url{http://www.gnu.org/software/octave/})
	\end{itemize}
    \end{enumerate}
\end{frame}
%%==================================================================F
\begin{frame}[label=vista_ppal]
 \frametitle{Vista principal de MATLAB}
  \begin{center}
   \includegraphics[width=0.7\paperwidth ]{images/matlab_vista_ppal}
  \end{center}
\end{frame}
%%==================================================================F 
\section{Variables y tipos de datos}
\begin{frame}[label=variables]
 \frametitle{Variables y tipos}
 \begin{enumerate}
  \item Todos los cálculos se guardan en variables
  \begin{itemize}
   \item \texttt{Identificador = Contenido}
   \begin{itemize}
    \item \texttt{a = 3}
    \item \texttt{b=a}
    \item \texttt{a = a + 1}
   \end{itemize}
  \end{itemize}
  \item Tipos básicos
  \begin{itemize}
   \item Enteros: \texttt{43, -1, 15}
   \item Reales: \texttt{8.5, .23, -2.7e3, 0.4e-2}
   \item Complejos: \texttt{3+2i, 9+1.4e-5i, 2+3j}
  \end{itemize}
 \end{enumerate}
\end{frame}
%%==================================================================F 
\subsection{Constantes}
\begin{frame}[label=constantes,shrink=5]
 \frametitle{Constantes}
 \begin{enumerate}
  \item Valores constantes predefinidos en MATLAB
  \begin{itemize}
   \item Infinito: \texttt{Inf, inf}
   \item ``Not a number'': \texttt{nan, NaN}
   \item \alert{¡No debería usarse como nombres de variables!}
  \end{itemize}
 \end{enumerate}
 \begin{center}
  \begin{scriptsize}
  \rowcolors{1}{ZurichBlue!20}{ZurichBlue!5}
   \begin{tabular}{|c|p{3.5cm}|c|}
   \multicolumn{1}{c}{\textbf{Nombre Constante}} 
   & \multicolumn{1}{c}{\textbf{Significado}} 
   & \multicolumn{1}{c}{\textbf{Valor}} \\ \toprule \hline
   \texttt{pi} & Constante $\pi$ & \texttt{3.14159265...} \\ \hline
   \texttt{e} & Constante $e$ & \texttt{2.71828182...} \\ \hline
   \texttt{i} & Unidad imaginaria & $\sqrt{-1}$ \\ \hline
   \texttt{j} & Unidad imaginaria & $\sqrt{-1}$ \\ \hline
   \texttt{intmax} & Entero más grande  & \texttt{2147483647} \\ \hline
   \texttt{intmin} & Entero más pequeño & \texttt{-2147483648} \\ \hline
   \texttt{realmax} & Número positivo más grande en doble precisión  &
\texttt{1.7977e+308} \\ \hline
   \texttt{realmin} & Número positivo más pequeño en doble precisión &
\texttt{2.2251e-308} \\ \hline
   \texttt{eps} & Precisión relativa en doble precisión & \texttt{2.2204e-16}
\\ \hline
  \end{tabular}
  \end{scriptsize}
 \end{center}
\end{frame}
%%==================================================================F 
\subsection{Enteros, reales, complejos y cadenas}
\begin{frame}[label=tipos_enteros]
 \frametitle{Tipos de datos enteros y reales}
 \begin{center}
  \begin{scriptsize}
  %\rowcolors{3}{ZurichBlue!20}{ZurichBlue!5}
   \begin{tabular}{|c|c|c|p{5cm}|}
   \multicolumn{1}{c}{ } & \multicolumn{1}{c}{\textbf{Tipo}} &
    \multicolumn{1}{c}{\textbf{Tamaño (bytes)}} &
    \multicolumn{1}{c}{\textbf{Rango}} \\ \toprule \hline
   \multirow{5}{*}{Enteros con signo} & \texttt{int8} & 1 & \texttt{-128 a 127}
    \\ \cline{2-4}
   & \texttt{int16} & 2 & \texttt{-32768 a 32767} \\ \cline{2-4}
   & \texttt{int32} & 4 & \texttt{-2147483648 a 2147483647} \\ \cline{2-4}
   & \multirow{2}{*}{\texttt{int64}} & \multirow{2}{*}{6} &
    \texttt{-9223372036854775808 a} \\ 
   & & & \texttt{9223372036854775807} \\ \hline
   \multirow{5}{*}{Enteros sin signo} & \texttt{uint8} & 1 & \texttt{0 a 255}
    \\ \cline{2-4}
   & \texttt{uint16} & 2 & \texttt{0 a 65535} \\ \cline{2-4}
   & \texttt{uint32} & 4 & \texttt{0 a 4294967295} \\ \cline{2-4}
   & \texttt{uint64} & 8 & \texttt{0 a 18446744073709551615} \\ \hline
   \multirow{2}{*}{Reales} & \texttt{single} & 4 & \texttt{-3.4028e+38 a
3.4028e+38} \\ \cline{2-4}
    & \texttt{double} & 8 & \texttt{-1.7977e+308 a 1.7977e+308} \\ \hline
  \end{tabular}
  \end{scriptsize}
 \end{center}
\end{frame}
%%==================================================================F 
\begin{frame}[label=tipo_reales]
 \frametitle{Formato de tipo de datos reales: \texttt{format}}
 \begin{enumerate}
  \item \texttt{format}: Especifica el formato de salida de los números.
  \item Representación del número $\pi$:
  \begin{itemize}
    \item \alert{\texttt{short}}: \texttt{3.1416}
    \begin{itemize}
     \item \texttt{short e}: \texttt{3.1416e+00}
     \item \alert{\texttt{short g}}: \\ \texttt{pi .\^\ [2; 4; 8; 16]\\
        9.8696\\97.409\\9488.5\\9.0032e+07}
    \end{itemize}
    \item \alert{\texttt{long}}: \texttt{3.14159265358979}
    \begin{itemize}
     \item \texttt{long e}: \texttt{3.14159265358979e+00}
     \item \alert{\texttt{long g}}: \texttt{3.14159265358979}                          
    \end{itemize}
  \end{itemize}
  \item \alert{\texttt{>> help format}}
 \end{enumerate}
\end{frame}
%%==================================================================F 
\begin{frame}[label=complejos]
 \frametitle{Números complejos}
 \begin{enumerate}
  \item Números complejos 
    \begin{itemize}
     \item Se usa la constante \alert{\texttt{i}} o \texttt{j} para expresar un número
        complejo
     \item Operaciones aritméticas:
      \begin{itemize}
        \item Suma: \texttt{z+w}
        \item Producto: \texttt{z*w}
        \item Cociente: \texttt{z/w}
        \item Conjugado: \texttt{z'}
      \end{itemize}
     \item Funciones:
      \begin{itemize}
       \item Parte real: \texttt{real(z)}
       \item Parte imaginaria: \texttt{imag(z)}
       \item Complejo conjugado: \texttt{conj(z)}
       \item Módulo: \texttt{abs(z)}
       \item Argumento: \texttt{angle(z)}
      \end{itemize}
    \end{itemize}
 \end{enumerate}
\end{frame}
%%==================================================================F
\subsection{Arrays: vectores y matrices}
\begin{frame}[label=arrays]
 \frametitle{¿Qué son los arrays?}
 \begin{enumerate}
  \item Generalmente son una colección de valores del mismo tipo
  \item Son un tipo de dato fundamental en Matlab
  \item Indexados mediante enteros o, incluso, un array de enteros
  \item Pueden ser filas o columnas
   \begin{itemize}
    \item \texttt{>> array = [0, 1, 2, 3]\\array = \\ \hspace*{0.5cm} 0 1 2 3}
    \item \texttt{>> column = [0; 1; 2; 3]\\column = \\ \hspace*{0.5cm} 0
        \\\hspace*{0.5cm} 1 \\\hspace*{0.5cm} 2 \\\hspace*{0.5cm} 3}
  \end{itemize}
\end{enumerate}
\end{frame}
%%==================================================================F
\begin{frame}[label=arrays2]
 \frametitle{Otras maneras de definir arrays}
 \begin{enumerate}
  \item Mediante una sucesión: \\\texttt{a=[0:4], a=(0:4), a=0:4}
  \item Con un incremento definido: \\\texttt{a=[0:2:10], a=(0:2:10), a=0:2:10}
  \item $N$ componentes equiespaciadas entre dos valores:
    \\\texttt{>> x=linspace(0,10,4) \\x = \\ \hspace*{0.5cm} 0.00000
\hspace*{0.25cm} 3.33333 \hspace*{0.25cm} 6.66667 \hspace*{0.25cm} 10.00000}
\end{enumerate}
\end{frame}
%%==================================================================F
\begin{frame}[label=matrices]
 \frametitle{Matrices}
 \begin{enumerate}
  \item Son vectores multi-dimensionales
  \item Para crear matrices, se crea dos vectores apilados en una columna
    \\\texttt{>> matriz = [1, 2, 3; 4, 5, 6] \\matriz = \\ \hspace*{0.5cm} 1
    \hspace*{0.25cm} 2 \hspace*{0.25cm} 3 \\ \hspace*{0.5cm} 4 \hspace*{0.25cm}
    5 \hspace*{0.25cm} 6}
 \end{enumerate}
\end{frame}
%%==================================================================F
\begin{frame}[label=cadenas]
 \frametitle{Cadenas de texto}
   \begin{enumerate}
    \item Se asignan con comillas simples: \texttt{s = '¡Hola, Mundo!'}
    \item Se pueden editar en el \emph{workspace}
   \end{enumerate}
  \begin{center}
   \includegraphics[height=3cm]{images/cadena_workspace}
  \end{center}
\end{frame}
%%==================================================================F
\section{Operadores}
%%==================================================================F 
\subsection{Asignación}
\begin{frame}[label=oper_asignacion]
 \frametitle{Operador asignación}
 \begin{enumerate}
 \item Operador asignación es el igual: \alert{\texttt{=}}
 \item Sirve para asignar valores a variables
  \begin{itemize}
   \item Ejemplo: \texttt{x=3}
   \item \alert{No confundir} con el operador de comparación:
        \alert{\texttt{==}} (un doble igual)
  \end{itemize}
 \item Si la variable no existe, se \alert{crea}
 \item Si existe, se \alert{sobre-escribe} y \alert{pierde} su valor anterior
 \end{enumerate}
\end{frame}
%%==================================================================F 
\subsection{Aritméticos}
\begin{frame}[label=oper_aritmeticos]
 \frametitle{Operadores aritméticos}
 \begin{center}
 \begin{small}
 \rowcolors{1}{ZurichBlue!20}{ZurichBlue!5}
  \begin{tabular}{|c|c|c|}
  \multicolumn{1}{c}{Operador} & \multicolumn{1}{c}{Símbolo} &
       \multicolumn{1}{c}{Ejemplo} \\ \toprule \hline
  Suma & \texttt{+} & \texttt{3+2} \\ \hline
  Resta & \texttt{-} & \texttt{3-2} \\ \hline
  Producto & \texttt{*} & \texttt{3*2} \\ \hline
  División derecha & \texttt{/} & \texttt{3/2} \\ \hline
  División izquierda & \texttt{\textbackslash} & \texttt{3\textbackslash 2} \\
\hline
  Potencia & \texttt{\^\ } & \texttt{3\^{}2} \\ \hline
  Resto de división entera & \texttt{rem(a,b)} & \texttt{rem(3,2)} \\
\hline
  \end{tabular}
 \end{small}
 \end{center}
\end{frame}
%%==================================================================F
\subsection{Relacionales y de igualdad}
\begin{frame}[label=oper_relacionales]
 \frametitle{Operadores relacionales y de igualdad}
 \begin{center}
 \begin{small}
\rowcolors{1}{ZurichBlue!20}{ZurichBlue!5}
  \begin{tabular}{|c|c|c|}
  \multicolumn{1}{c}{Operador} & \multicolumn{1}{c}{Símbolo} &
       \multicolumn{1}{c}{Ejemplo} \\ \toprule \hline
  Menor que & \texttt{<} & \texttt{3<2} \\ \hline
  Mayor que & \texttt{>} & \texttt{3>2} \\ \hline
  Menor o igual que & \texttt{<=} & \texttt{3<=2} \\ \hline
  Mayor o igual que & \texttt{>=} & \texttt{3>=2} \\ \hline
  Es igual & \texttt{==} & \texttt{3==2} \\
\hline
  Es diferente & \texttt{$\sim$=} & \texttt{3$\sim$=2} \\ \hline 
  \end{tabular}
 \end{small}
 \end{center}
 Devuelven un \alert{booleano} (True o False)
\end{frame}
%%==================================================================F
\subsection{Lógicos}
\begin{frame}[label=oper_logicos]
 \frametitle{Operadores lógicos}
\begin{columns} % contents are top vertically aligned
 \begin{column}{3cm} % each column can also be its own environment
   \begin{enumerate}[<+->]%
    \item \alert<1>{Y lógico}
    \item \alert<2>{O lógico}
    \item \alert<3>{XOR}
    \item \alert<4>{NO}
   \end{enumerate}
 \end{column}
 \begin{column}{7cm} % alternative top-align that's better for graphics
   \only<1>{
    \begin{center}
    \begin{small}
    \rowcolors{1}{ZurichBlue!20}{ZurichBlue!5}
    \begin{tabular}{|p{2cm}|p{2cm}|p{2cm}|}
    \multicolumn{3}{c}{\texttt{\textbf{C = A \& B}}} \\ \toprule \hline
   \textbf{A} & \textbf{B} & \textbf{C} \\ \hline
    1 & 1 & 1 \\ \hline
    1 & 0 & 0 \\ \hline
    0 & 1 & 0 \\ \hline
    0 & 0 & 0 \\ \hline
    \end{tabular}
    \end{small}
    \end{center}}

   \only<2>{
    \begin{center}
    \begin{small}
    \rowcolors{1}{ZurichBlue!20}{ZurichBlue!5}
    \begin{tabular}{|p{2cm}|p{2cm}|p{2cm}|}
    \multicolumn{3}{c}{\texttt{\textbf{C = A | B}}} \\ \toprule \hline
   \textbf{A} & \textbf{B} & \textbf{C} \\ \hline
    1 & 1 & 1 \\ \hline
    1 & 0 & 1 \\ \hline
    0 & 1 & 1 \\ \hline
    0 & 0 & 0 \\ \hline
    \end{tabular}
    \end{small}
    \end{center}}

   \only<3>{
    \begin{center}
    \begin{small}
    \rowcolors{1}{ZurichBlue!20}{ZurichBlue!5}
    \begin{tabular}{|p{2cm}|p{2cm}|p{2cm}|}
    \multicolumn{3}{c}{\texttt{\textbf{C = xor(A,B)}}} \\ \toprule \hline
   \textbf{A} & \textbf{B} & \textbf{C} \\ \hline
    1 & 1 & 0 \\ \hline
    1 & 0 & 1 \\ \hline
    0 & 1 & 1 \\ \hline
    0 & 0 & 0 \\ \hline
    \end{tabular}
    \end{small}
    \end{center}}

   \only<4>{
    \begin{center}
    \begin{small}
    \rowcolors{1}{ZurichBlue!20}{ZurichBlue!5}
    \begin{tabular}{|p{3cm}|p{3cm}|}
    \multicolumn{2}{c}{\texttt{\textbf{C =$\sim$ A}}} \\ \toprule \hline
   \textbf{A} & \textbf{C} \\ \hline
    1 & 0 \\ \hline
    0 & 1 \\ \hline
    \end{tabular}
    \end{small}
    \end{center}}

 \end{column}
\end{columns}
\end{frame}
%%==================================================================F
\begin{frame}[label=oper_prioridad,shrink=5]
 \frametitle{Prioridad entre operadores}
 \begin{center}
 \begin{scriptsize}
 \rowcolors{1}{ZurichBlue!20}{ZurichBlue!5}
  \begin{tabular}{|l|c|l|}
  \multicolumn{1}{c}{} & \multicolumn{1}{c}{} &
       \multicolumn{1}{c}{\textbf{A igualdad prioridad...}} \\
  \toprule \hline
  Paréntesis, XOR & \texttt{(), xor()} & \\ \hline
  No lógico & \texttt{$\sim$} & Derecha a izquierda \\ \hline
  Potencia & \texttt{\^\ } & Izquierda a derecha \\ \hline
  Multiplicación, división & \texttt{*,/} & Izquierda a derecha \\ \hline
  Suma, resta, cambio de signo & \texttt{+,-} & Izquierda a derecha \\ \hline
  Relacionales & \texttt{<, <=, >, >=} & Izquierda a derecha \\ \hline
  Igualdad & \texttt{==, $\sim$=} & Izquierda a derecha \\ \hline
  Y lógico & \texttt{\&} & Izquierda a derecha \\ \hline
  O lógico & \texttt{|} & Izquierda a derecha \\ \hline
  Asignación & \texttt{=} & Izquierda a derecha \\ \hline
  \end{tabular}
 \end{scriptsize}
 \end{center}
    Es mucho mejor hacer explícito el orden de las operaciones usando el
    paréntesis: 
  \begin{itemize}
    \item \texttt{1 + 2 * 3 / 4 - 5}
    \item \texttt{(1 + (2 * 3) / 4) - 5}            
  \end{itemize}
\end{frame}
%%==================================================================F
\subsection{Matrices}
\begin{frame}[label=array_operaciones]
 \frametitle{Matrices}
 \begin{enumerate}
  \item Suma y resta:
    \begin{itemize}
     \item \texttt{A+B}: Operación matricial elemento a elemento
     \item \texttt{A+c}: Operación escalar. (Suma $c$ a cada elemento)
    \end{itemize}
  \item Producto:
    \begin{itemize}
     \item \texttt{A*B}: Operación matricial elemento a elemento
     \item \texttt{A*c}: Operación escalar. (Multiplica $c$ a cada elemento)
    \end{itemize}
  \item Producto \alert{elemento a elemento}: \texttt{A.*B} $\Rightarrow
    A_{ij}\cdot B_{ij}$
  \item Potenciación: \texttt{A\^{}3} $\Rightarrow A^3=A\cdot A\cdot A$
    (Producto matricial)
  \item Potenciación \alert{elemento a elemento}:
    \begin{itemize}
     \item \texttt{A.\^{}3}: $\Rightarrow A_{ij}^3$
     \item \texttt{A.\^{}B}: $\Rightarrow A_{ij}^{B_{ij}}$
    \end{itemize}
 \end{enumerate}
\end{frame}
%%==================================================================F
\begin{frame}[label=array_operaciones2]
 \frametitle{Matrices}
 \begin{enumerate}
  \item División:
    \begin{itemize}
     \item \texttt{A/c} Operación escalar. Divide entre $c$ cada elemento
     \item \texttt{c/A} \alert{ERROR!}
     \item \texttt{A/B} $\Leftrightarrow$ \texttt{A*inv(B)}
     \item \texttt{A\textbackslash B} $\Leftrightarrow$ \texttt{inv(A)*B}
    \end{itemize}
  \item División elemento a elemento:
    \begin{itemize}
     \item \texttt{A./B} $\Rightarrow A_{ij}/B_{ij}$
     \item \texttt{c./A} $\Rightarrow c/A_{ij}$
    \end{itemize}
  \item Traspuesta: \texttt{A'}
 \end{enumerate}
\end{frame}
%%==================================================================F 
\section{Programación en MATLAB}
\begin{frame}[label=programacion]
\frametitle{Programación}
 \begin{enumerate}
  \item Secuencia de comandos que se almacenan en un archivo
  \item Se usa un editor de código (Matlab)
  \item El programa se guardaría como \alert{\texttt{programa.m}} en el 'pwd'
  \item El programa se ejecuta con \\\texttt{>> programa}
  \item Los comentarios empiezan con el signo \texttt{
      \%:\\\textcolor{Comentario}{\% Esto es un comentario}}
 \end{enumerate}
\end{frame}
%%==================================================================F 
\subsection{Bifurcaciones}
\begin{frame}[label=bifuraciones,fragile]
 \frametitle{Bifurcaciones}
 \begin{itemize}
    \item Ejecutar una parte de un programa si se cumple una \alert{condición}.
    \item O ejecutar otra si no se cumple.
 \end{itemize}
 \defverbatim[colored]\lstI{
   \begin{lstlisting}[style=mio]
       if (condicion)
        sentencias % Condicion verdadera
       else
        sentencias % Condicion falsa
       end
   \end{lstlisting}
 }
 \begin{small}
  \lstI
 \end{small}
 \begin{itemize}
    \item Cuando no hay que hacer nada si la expresión booleana es falsa, el
      \alert{\texttt{else}} se puede obviar.
 \end{itemize}
\end{frame}
%%==================================================================F
\begin{frame}[label=bifuraciones_multiples,fragile]
 \frametitle{Bifurcaciones múltiples}
 \begin{itemize}
    \item A veces se hace necesario considerar varias condiciones excluyentes,
    \item y no tiene sentido crear un \alert{\texttt{if}} para cada condición.
 \end{itemize}
 \defverbatim[colored]\lstI{
   \begin{lstlisting}[style=mio]
      if (condicion 1)
        sentencias
      elseif (condicion 2)
        sentencias
      else
        sentencias
      end
   \end{lstlisting}
 }
 \begin{small}
  \lstI
 \end{small}
\end{frame}
%%==================================================================F
\begin{frame}[label=case,fragile]
 \frametitle{Bifurcaciones múltiples}
 \begin{itemize}
    \item Cuando tenemos muchas condiciones posibles una opción puede ser
      \alert{\texttt{switch-case}}.
 \end{itemize}
 \defverbatim[colored]\lstI{
   \begin{lstlisting}[style=mio]
      switch variable
        case 'valor 1'
          sentencias
        case 'valor 2'
          sentencias
        ...
        otherwise
          sentencias
      end
 \end{lstlisting}
 }
 \begin{small}
   \lstI
 \end{small}
\end{frame}
%%==================================================================F
\subsection{Iteraciones}
\begin{frame}[label=iteraciones]
 \frametitle{Iteraciones}
 \begin{enumerate}
  \item En muchas ocasiones querremos repetir las mismas instrucciones varias
    veces (\alert{bucle}).
  \item Un bucle \alert{\texttt{for}} itera sobre una lista (vector), y repite
    los comandos un número determinado de veces.
  \item Un bucle \alert{\texttt{while}} repite los comandos mientras se cumpla
    una condición.
  \item Cualquier bucle \alert{\texttt{for}} se puede transformar en un bucle
    \alert{\texttt{while}}.
 \end{enumerate}
\end{frame}
%%==================================================================F
\begin{frame}[label=bucle_for,fragile]
 \frametitle{Bucle \texttt{for}}
 \begin{itemize}
  \item \alert{\texttt{for}} Ejecuta el código un número de veces utilizando un
    iterador o índice.
 \defverbatim[colored]\lstI{
  \begin{lstlisting}[style=mio]
    for indice = valor_inicial:incremento:valor_final
      sentencias
     end
 \end{lstlisting}
 }
 \begin{small}
   \lstI
 \end{small}
  \item El índice se inicializa a \texttt{valor\_inicial} y se incrementa un
    valor \texttt{incremento} hasta alcanzar el \texttt{valor\_final}
  \item Si se omite \texttt{incremento} se asume que su valor es \texttt{1}.
 \end{itemize}
\end{frame}
%%==================================================================F
\begin{frame}[label=bucle_while,fragile]
 \frametitle{Bucle \texttt{while}}
 \begin{itemize}
  \item \alert{\texttt{while}} Ejecuta el código hasta que una cierta condición
    es \alert{falsa}.
 %\begin{small}
 \begin{verbatim}
    while condicion
      sentencias
    end
 \end{verbatim}
 %\end{small}
  \item \alert{¡OJO!} Cambiar la condición dentro del bucle puede causar
    \alert{bucles infinitos}.
 \end{itemize}
\end{frame}
%%==================================================================F
\begin{frame}[label=break_continue]
 \frametitle{Interrupción de iteraciones: \texttt{break} y \texttt{continue}}
 \begin{itemize}
  \item Puede ser necesario \alert{salir} de un bucle de manera abrupta.
  \begin{itemize}
   \item Podemos usar la instrucción \alert{\texttt{break}}.
  \end{itemize}
  \item En cualquier iteración, podemos \alert{“saltar”} directamente a la siguiente.
  \begin{itemize}
   \item Se usa la instrucción \alert{\texttt{continue}}.
  \end{itemize}
  \item Es también posible salir de una función o programa completo con
    \alert{\texttt{return}}.
  \item Siempre es mejor diseñar el bucle para que la salida sea una
     consecuencia lógica del bucle.
 \end{itemize}
\end{frame}
%%==================================================================F
\begin{frame}[label=break,fragile]
 \frametitle{Ejemplo de \texttt{break}}
 \begin{itemize}
  \item \alert{\texttt{break}} obliga a salir del bucle y continuar con la
    ejecución del programa.
   \defverbatim[colored]\lstI{
     \begin{lstlisting}[style=mio]
      ...
      for i=1:100
        if (i==11)
          break;     % Finaliza el bucle
        end
        disp(i);
      end
      ...
    \end{lstlisting}
  }
 \begin{small}
  \lstI
 \end{small}
  \item El valor de i es \alert{\texttt{i=11}}.
 \end{itemize}
\end{frame}
%%==================================================================F
\begin{frame}[label=continue,fragile]
 \frametitle{Ejemplo de \texttt{continue}}
 \begin{itemize}
  \item \alert{\texttt{continue}} obliga la finalización de esa repetición
del bucle y comienza con la siguiente.
 \begin{scriptsize}
   \defverbatim[colored]\lstI{
     \begin{lstlisting}[style=mio]
      ...
      for i=1:100
        if (i==11)
          continue;     % Continua con la siguiente iteracion
        end
        disp(i);
      end
      ...
    \end{lstlisting}
  }
  \lstI
 \end{scriptsize}
  \item El valor de i es \alert{\texttt{i=100}}.
 \end{itemize}
\end{frame}
%%==================================================================F
\section{Funciones}
%%==================================================================F
\subsection{Definidas por el usuario}
\begin{frame}[label=func_usuario]
 \frametitle{Funciones definidas por el usuario}
 \begin{itemize}
  \item Hay que manejar tres conceptos básicos:
  \begin{enumerate}
   \item Argumentros de entrada.
   \item Variables locales.
   \item Argumentos de salida.
  \end{enumerate}
  \item Oculta los detalles del cálculo.
  \item Facilita el diseño (concept \alert{top-down}).
  \item Facilita entender un programa $\Rightarrow$ Más fácil de modificar en
    el futuro o por otras personas.
 \end{itemize}
\end{frame}
%%==================================================================F
\begin{frame}[label=func_declaracion,fragile]
 \frametitle{Declaración de una función}
 \begin{itemize}
  \item Para declarar una variable se utiliza la palabra clave    
    \alert{\texttt{function}}
  \item Por norma general, acabaremos la función con la palabra clave
    \alert{\texttt{end}}
  \item Es una buena praxis el comentar la función con un texto claro, conciso e
    informativo
 \end{itemize}
   \defverbatim[colored]\lstI{
     \begin{lstlisting}[style=mio]
     function s = MiFuncion (E)
        % La documentacion de las funciones se incluye
        % justo despues de la cabecera.
        % E: Argumentos de Entrada.
        % s: Argumentos de salida.
        ...
     end
    \end{lstlisting}
  }
 \begin{small}
  \lstI
 \end{small}
\end{frame}
%%==================================================================F
\begin{frame}[label=func_ficheros]
 \frametitle{Funciones en ficheros}
 \begin{itemize}
  \item Cada función se guarda en un fichero '\texttt{.m}'.
  \item El nombre del fichero es igual que la función: \texttt{MiFuncion.m}
  \item Si sólo hay una función dentro del fichero, se puede omitir
    \alert{\texttt{end}} al final de la función.
  \item Si hay varias funciones, sólo la primera queda visible fuera del
    fichero.
  \item La primera línea del fichero que no sea un comentario tiene que ser la
    declaración de la función.
  \item Para llamar a una función, tenemos que estar en el mismo directorio que
    su fichero \texttt{funcion.m}.
 \end{itemize}
\end{frame}
%%==================================================================F
\begin{frame}[label=func_variables,shrink=1]
 \frametitle{Variables locales y globales}
 \begin{itemize}
  \item Las variables \alert{locales} están definidas dentro del cuerpo de una
    función y no son visibles fuera de él.
  \item Las variables \alert{globales} son las que definimos fuera del cuerpo de
    una función.
  \item Si hay una variable local con el mismo nombre que una variable global,
    la variable global ya no está disponible dentro de esa función.
  \item Podemos crear una variable global desde dentro del cuerpo de una función
    usando la palabra clave global: \\\hspace{15mm}\texttt{global m; m = 1;}
  \item Es muy \alert{poco recomendable} crear variables globales dentro de
    funciones.
  \begin{itemize}
    \item Lo que pasa en la función, queda dentro de la función \texttt{;-)}
  \end{itemize}
 \end{itemize}
\end{frame}
%%==================================================================F
\begin{frame}[label=func_argumentos]
 \frametitle{Argumentos de entrada y salida}
 \begin{itemize}
  \item Podemos utilizar tantos argumentos de entrada y salida como queramos.
  \item Podemos especificar que algunos de los argumentos de nuestras funciones
    sean opcionales.
  \item Para saber cuántos argumentos se han pasado usamos   
    \alert{\texttt{nargin}}.
  \item Para saber cuántos argumentos de salida nos piden, usamos
    \alert{\texttt{nargout}}.
 \end{itemize}
\end{frame}
%%==================================================================F
\subsection{Funciones de la librería de MATLAB}
\begin{frame}[label=func_builtin,shrink=10]
 \frametitle{Funciones matemáticas}
 \begin{enumerate}
  \item Cáculos básicos
   \begin{itemize}
    \item Valor absoluto: \texttt{abs(x)}
    %\item Resto de la división: \texttt{rem(x,y)}
    \item Raíz cuadrada: \alert{\texttt{sqrt(x)}}
    \item Logaritmo neperiano: \texttt{log(x)}
    \item Función exponencial: \alert{\texttt{exp(x)}}
    \item Mayor entero menor o igual a x: \texttt{floor(x)}
    \item Menor entero mayor o igual a x: \texttt{ceil(x)}
    \item Redondea al entero más cercano: \texttt{round(x)}
   \end{itemize}
  \item Cálculo trigonométrico
   \begin{itemize}
    \item Seno y arcoseno: \texttt{\alert{sin(x)}, asin(x)}
    \item Coseno y arcocoseno: \texttt{\alert{cos(x)}, acos(x)}
    \item Tangente y arcotangente: \texttt{\alert{tan(x)}, atan(x),
      \alert{atan2(x,y)}}
    %\item Cotangente y arcocotangente: \texttt{cot(x), acot(x)}
    %\item Secante y arcosecante: \texttt{csc(x), acsc(x)}
    %\item Cosecante y arcocosecante: \texttt{sec(x), asec(x)}
 \end{itemize}
  \item Cálculo estadístico
   \begin{itemize}
    \item Media: \alert{\texttt{mean(x)}}
    \item Desviación típica: \alert{\texttt{std(x)}}
 \end{itemize}
\end{enumerate}
\end{frame}
%%==================================================================F
\subsubsection{Matrices}
\begin{frame}[label=array_funciones,shrink=10]
 \frametitle{Matrices}
 \begin{itemize}
  \item Producto escalar: \texttt{dot(u,v)}
  \item Producto vectorial: \texttt{cross(u,v)}
  \item Tamaño de un array o matriz: \texttt{length(v)}
  \item Suma de componentes: \texttt{sum(v)}
  \item Módulo: \texttt{norm(v)}
  \item Inversa de una matriz: \texttt{Inv(A)}
  \item Determinante: \texttt{det(A)}
  \item Traza: \texttt{trace(A)}
  \item Matriz identidad $I_n$: \alert{\texttt{eye(n)}}
  \item Matriz de ceros: \texttt{zeros(m,n), zeros(n)}
  \item Matriz de unos: \texttt{ones(m,n), ones(n)}
  \item Generar matriz diagonal: \texttt{diag(v)}
 \end{itemize}
\end{frame}
%%==================================================================F
\begin{frame}[label=func_builtin2,shrink=5]
 \frametitle{Descomposición de matrices}
 \begin{enumerate}
   \item Descomposición \alert{LU}: Descompone la matriz A en el producto
      $A=L\cdot U$, siendo $L$ una matriz triangular inferior y $U$ una
      triangular superior.
      \\\hspace{1cm}\texttt{[L,U]=lu(A)}
   \item Factorización de \alert{Cholesky}: Devuelve la matriz triangular
      superior $R$ tal que $R'\cdot R=A$, siempre que A sea definida positiva.
      \\\hspace{1cm}\texttt{R=chol(A)}
   \item Descomposición \alert{ortogonal}: Devuelve la matriz triangular
      superior $R$ de igual dimensión que $A$ y la matriz ortogonal $Q$ tal que
      $A=Q\cdot R$.
      \\\hspace{1cm}\texttt{[Q,R]=qr(A)}
 \end{enumerate}
\end{frame}
%%==================================================================F
\section{Entrada/Salida}
\begin{frame}[label=io]
 \frametitle{Entrada y salida de datos}
 \begin{enumerate}
   \item La entrada y salida más comunes son:
   \begin{itemize}
    \item Estándar.
    \item De/a un fichero.
   \end{itemize}
   \item Se pueden introducir:
   \begin{itemize}
    \item Números.
    \item Cadenas.
    \item Vectores.
    \item Matrices.
   \end{itemize}
 \end{enumerate}
\end{frame}
%%==================================================================F
\subsection{E/S estándar}
\begin{frame}[label=io_std]
 \frametitle{Entrada y salida estándar}
 \begin{enumerate}
  \item Entrada:
  \begin{itemize}
   \item \texttt{variable = input('Mensaje para el usario')}
   \item El resultado se procesa al presionar \alert{$<$enter$>$}
   \item Para recoger una cadena: 
   \begin{itemize}
    \item \texttt{variable = input('Mensaje para el usario','s')}
    \item Podemos pedir al usuario que escriba el valor entre comillas simples
   \end{itemize}
  \end{itemize}
  \item Salida:
  \begin{itemize}
   \item Solución simple: \alert{no pongas} '\alert{\texttt{;}}'
   \item Para enviar mensajes:
   \begin{itemize}
     \item \texttt{disp('Mensaje para el usario')}
   \end{itemize}
   \item Para mostrar variables:
   \begin{itemize}
     \item \texttt{disp(variable)}
   \end{itemize}
  \end{itemize}
 \end{enumerate}
\end{frame}
%%==================================================================F
\subsection{E/S con formato}
\begin{frame}[label=salida_std]
 \frametitle{Salida y entrada con formato}
 \begin{enumerate}
    \item Salida con formato:
  \begin{itemize}
    \item \alert{\texttt{printf(TEMPLATE, variables)}}: Salida estándar.
    \item \alert{\texttt{sprintf}}: La salida es una cadena de caracteres.
    \item \alert{\texttt{fprintf}$^*$}: La salida se realiza a un fichero
  \end{itemize}
  \item Entrada con formato:
  \begin{itemize}
    \item \alert{\texttt{scanf}}: Entrada estándar (no muy útil).
    \item \alert{\texttt{sscanf}}: La entrada es una cadena de caracteres.
    \item \alert{\texttt{fscanf}$^*$}: Toma la entrada de un fichero.
  \end{itemize}
  \item \alert{$^{*}$}Necesitan las funciones:
  \begin{itemize}
    \item Abrir archivo: \alert{\texttt{fopen}}
    \item Cerrar archivo: \alert{\texttt{fclose}}
  \end{itemize}
 \end{enumerate}
\end{frame}
%%==================================================================F
\begin{frame}[label=template]
 \frametitle{\texttt{TEMPLATE}: Dando formato}
 \begin{enumerate}
  \item Formato de datos:
  \begin{itemize}
    \item Enteros, lógicos, reales con decimales cero: \alert{'\texttt{\%d}'}
    \item Cualquier dato numérico con 6 decimales: \alert{'\texttt{\%f}'}
    \item Escribe cadenas: \alert{'\texttt{\%s}'}
  \end{itemize}
  \item Caracteres de control:
  \begin{itemize}
    \item Salto de línea: \alert{\texttt{'\textbackslash n'}}
    \item Retorno de carro: \texttt{'\textbackslash r'}
    \item Tabulación: \alert{\texttt{'\textbackslash t'}}
    \item Espacio hacia atrás: \texttt{'\textbackslash b'}
  \end{itemize}
 \end{enumerate}
\end{frame}
%%==================================================================F
\subsection{Importar y Exportar}
\begin{frame}[label=io_otras]
 \frametitle{Otras salidas y entradas de datos}
 \begin{enumerate}
  \item Importar y exportar datos:
  \begin{itemize}
    \item Cargar datos formato \alert{Matlab} o en fichero \alert{ASCII}
      delimitado por espacios: \alert{\texttt{load}}
    \item Guardar datos formato \alert{Matlab} o en fichero \alert{ASCII}
      delimitado por espacios: \alert{\texttt{save}}
  \end{itemize}
 \end{enumerate}
\end{frame}
%%==================================================================F
\againframe{portada}
%%==================================================================F

%\section{Resumen}
